% Source: https://artofproblemsolving.com/community/c3h2684047_a_bunch_of_problems
% Stable identity: aops-topic-2684047-q01
Let the radius be $r$, the centre $O$ and the other vertices $A,B$. The triangle has side $r$, so $\angle AOB=60^\circ$. The three regions are the triangle $OAB$, the circular segment between chord $\overline{AB}$ and the minor arc $AB$, and the rest of the disc. Their areas are
\[[OAB]=\dfrac{\sqrt3}{4}r^2\approx0.433r^2,\qquad \text{segment}=\dfrac{60}{360}\pi r^2-\dfrac{\sqrt3}{4}r^2\approx0.0906r^2,\qquad\text{rest}=\dfrac56\pi r^2.\]
The segment is the smallest, so the shaded fraction of the disc is
\[\dfrac{\dfrac16\pi r^2-\dfrac{\sqrt3}{4}r^2}{\pi r^2}=\dfrac16-\dfrac{\sqrt3}{4\pi},\]
which is $\boxed{\left(\dfrac{50}{3}-\dfrac{25\sqrt3}{\pi}\right)\%\approx2.88\%}$.

\begin{asy}
usepackage("amsmath");
size(3.5cm);
defaultpen(fontsize(10pt));
pair O=(0,0),A=dir(30),B=dir(-30);
fill(arc(O,1,-30,30)--cycle,mediumgray);
draw(unitcircle,linewidth(0.8));draw(O--A--B--cycle,linewidth(0.8));
dot(O);label("$O$",O,W);label("$A$",A,NE);label("$B$",B,SE);
label("$r$",(O+A)/2,NW,fontsize(8pt));label("$r$",(O+B)/2,SW,fontsize(8pt));
draw(arc(O,0.2,-30,30));label("$60^\circ$",(0.22,0),E,fontsize(7pt));
\end{asy}
